\documentclass[main.tex]{subfiles}


\begin{document}

%from pagenumber to up
% \phantomsection 
% \hypertarget{toc}{}

 

 

 
   

\section{Изопроцессы}


\textbf{Процесс}~--- это переход газа\footnote{Более точно, макроскопического тела.} из одного состояния в другое. Состояние газа определяется его количеством вещества~($\nu$), давлением~($P$), объёмом~($V$) и температурой~($T$); эти параметры связаны друг с другом уравнением Менделеева"--~Клапейрона. В дальнейшем подразумевается, что процессы проводят с одним и тем же газом (то есть $\nu=\mathrm{const}$).

%При описании природных явлений и процессов в устройствах часто имеют дело с газом, у которого \emph{масса}~($m$) \emph{и молярная масса}~($M$) \emph{остаются неизменными}: в таком случае состояние газа определяется \emph{тремя} макроскопическими параметрами: давлением~($P$), объёмом~($V$) и температурой~($T$). В дальнейшем подразумевается, что процессы проводят с одним и тем же газом.

% При описании природных явлений и процессов в технических устройствах имеют дело с самыми разнообразными процессами, происходящими с газами.

\textbf{Изопроцесс}~--- это процесс, в котором значение одного из макропараметров ($P$, $V$ или $T$) остается неизменным. На рис.\,\ref{fig:p71} изображены три сосуда с одинаковыми газами, у которых поддерживаются постоянными разные параметры.

\begin{figure}[H]
    \centering

    \begin{tikzpicture}[font=\small,            line cap=round,                             >={Stealth[inset=0pt,round,length=3mm, angle'=20]},vec/.style={>={Stealth[inset=0pt,round,length=3.5mm, angle'=20]}}]


\filldraw[lightgray,semitransparent] (0,0) rectangle++ (11,-.3);  
\draw (0,0) --++ (11,0);

% T=const

%plunger
\fill[lightgray] (1,1.5) rectangle++ (0.9,.15);
\filldraw[lightgray,draw=black] (1,1.5) --++ (0.9,0) ++ (0,.15) --++ (-.9,0);

\fill[lightgray] (3,1.5) rectangle++ (-0.9,.15);
\filldraw[lightgray,draw=black] (3,1.5) --++ (-0.9,0) ++ (0,.15) --++ (.9,0);


%temp
\path[] (1.9,1.4) arc [start angle=180, end angle=240, radius=.1] coordinate (T) --++ (0,-.1) arc [start angle=180, end angle=360, radius={cos(60)*.1}] --++ (0,.1) arc [start angle=300, end angle=360, radius=.1] --++ (0,1) arc [start angle=0, end angle=180, radius=.1] -- cycle;

\fill[red] (T) --++ (0,-.1) arc [start angle=180, end angle=360, radius={cos(60)*.1}] --++ (0,.1) arc [start angle=0, end angle=180, radius={cos(60)*.1}];
\draw[] (1.9,1.4) arc [start angle=180, end angle=240, radius=.1] coordinate (T) --++ (0,-.1) arc [start angle=180, end angle=360, radius={cos(60)*.1}] --++ (0,.1) arc [start angle=300, end angle=360, radius=.1] --++ (0,1) arc [start angle=0, end angle=180, radius=.1] -- cycle;
\draw[red] (T) ++ ({cos(60)*.1},0) --++ (0,.5) coordinate (Tr);
\draw[lightgray] (Tr) --++ (0,.5);

%gas
\draw (1,0) ++ (0,2) --++ (0,-2) --++ (2,0) --++ (0,2);

\begin{scope}
    \clip[] (1.9,1.5) --++ (0,-.1) arc [start angle=180, end angle=240, radius=.1] coordinate (T) --++ (0,-.1) arc [start angle=180, end angle=360, radius={cos(60)*.1}] --++ (0,.1) arc [start angle=300, end angle=360, radius=.1] --++ (0,.1) --++ (.9,0) --++ (0,-1.5) --++ (-2,0) --++ (0,1.5) -- cycle;
\foreach \i in {1,...,100}
  \fill[Green,shift={(1.04,0.04)}] ({rnd*1.92cm}, {rnd*1.42cm}) circle (.02);
\end{scope}


\node[fill=white] at (2,.75) {$T=\mathrm{const}$};  

%%%%%%%%%%%%%%%%%%%%%%%%%%%%%%%%%%%%%%%%
%%%%%%%%%%%%%%%%%%%%%%%%%%%%%%%%%%%%%%%%

%V=const

\begin{scope}[xshift=3.5cm]

%limit
\fill[gray] (1,1.5) ++ (0,-.2) --++ (.2,.2) --++ (-.2,0) --++ (0,-.2); 
\draw (1,1.5) ++ (0,-.2) --++ (.2,.2);
\fill[gray] (1,1.65) ++ (0,.2) --++ (.2,-.2) --++ (-.2,0) --++ (0,.2); 
\draw (1,1.65) ++ (0,.2) --++ (.2,-.2);

\fill[gray] (3,1.5) ++ (0,-.2) --++ (-.2,.2) --++ (.2,0) --++ (0,-.2); 
\draw (3,1.5) ++ (0,-.2) --++ (-.2,.2);
\fill[gray] (3,1.65) ++ (0,.2) --++ (-.2,-.2) --++ (.2,0) --++ (0,.2); 
\draw (3,1.65) ++ (0,.2) --++ (-.2,-.2);


%plunger
\fill[lightgray] (1,1.5) rectangle++ (2,.15);
\filldraw[lightgray,draw=black] (1,1.5) --++ (2,0) ++ (0,.15) --++ (-2,0);

%gas
\draw (1,0) ++ (0,2) --++ (0,-2) --++ (2,0) --++ (0,2);

\foreach \i in {1,...,100}
  \fill[Green,shift={(1.04,0.04)}] ({rnd*1.92cm}, {rnd*1.42cm}) circle (.02);


\node[fill=white] at (2,.75) {$V=\mathrm{const}$};  
    
\end{scope}


%%%%%%%%%%%%%%%%%%%%%%%%%%%%%%%%%%%%%%%%
%%%%%%%%%%%%%%%%%%%%%%%%%%%%%%%%%%%%%%%%

%P=const

\begin{scope}[xshift=7cm]

%load
\filldraw[black,rounded corners=1pt] (2,1.65) --++ (-.2,0) --++ (0,.5) --++ (.15,0) --++ (0,.1) --++ (-.1,0) --++ (0,.1) --++ (.3,0) --++ (0,-.1) --++ (-.1,0) --++ (0,-.1) --++ (.15,0) --++ (0,-.5) -- cycle;


%plunger
\fill[lightgray] (1,1.5) rectangle++ (2,.15);
\filldraw[lightgray,draw=black] (1,1.5) --++ (2,0) ++ (0,.15) --++ (-2,0);

%gas
\draw (1,0) ++ (0,2) --++ (0,-2) --++ (2,0) --++ (0,2);

\foreach \i in {1,...,100}
  \fill[Green,shift={(1.04,0.04)}] ({rnd*1.92cm}, {rnd*1.42cm}) circle (.02);


\node[fill=white] at (2,.75) {$P=\mathrm{const}$};  
    
\end{scope}


    \end{tikzpicture}
\caption{Газы с разными фиксированными параметрами}
\label{fig:p71}
\end{figure}

Трем случаям, проиллюстрированным на рисунке, соответствуют три вида изопроцессов.
\begin{enumerate}
    \item \textbf{Изотермический} процесс (рис.\,\ref{fig:p71}, слева) идет при постоянной температуре газа: $T=\mathrm{const}$. График\footnote{Графики зависимости одного параметра от другого в изопроцессах получены с использованием уравнения Менделеева"--~Клапейрона.} такого процесса называется \emph{изотермой}. На рис.\,\ref{fig:p72} (слева) изображены две изотермы газа с разными температурами на $PV$"=диаграмме: \emph{чем выше температура, тем дальше от начала~$O$ системы лежит соответствующая изотерма}. 
    \item \textbf{Изохорный} процесс (рис.\,\ref{fig:p71}, посередине) идет при постоянном объеме газа: $V = \mathrm{const}$. График такого процесса~--- \emph{изохора}. На рис.\,\ref{fig:p72} (посередине) изображены две изохоры газа с разными объемами на $PT$"=диаграмме: \emph{чем больше объем, тем ближе к оси~$T$ лежит соответствующая изохора}. 
    \item \textbf{Изобарный} процесс (рис.\,\ref{fig:p71}, справа) идет при постоянном давлении газа: $P = \mathrm{const}$. График такого процесса есть \emph{изобара}. На рис.\,\ref{fig:p72} (справа) изображены две изобары газа с разными давлениями на $VT$"=диаграмме: \emph{чем больше давление, тем ближе к оси~$T$ лежит соответствующая изобара}.
\end{enumerate}

% Графики процессов обычно изображают в $PV$-, $PT$- и $VT$"=системах (рис.\,\ref{fig:p72}). 

\begin{figure}[H]
    \centering

    \begin{tikzpicture}[font=\small,            line cap=round,                             >={Stealth[inset=0pt,round,length=3mm, angle'=20]},vec/.style={>={Stealth[inset=0pt,round,length=3.5mm, angle'=20]}}]


%изотерма
\draw[->] (-.3,0) --++ (0:3cm) node[below,black] {$V$};
\draw[->] (0,-.3) --++ (90:3cm) node[left,black] {$P$};
\node[below left] at (0,0) {$O$};

\draw[red,thick,domain=.35:1.75,samples=1000] plot (\x,{1/(2*\x)}) node[right,black] {$T_1$};
\draw[red,thick,domain=.65:1.75,samples=1000] plot (\x,{1.5/\x}) node[right,black] {$T_2>T_1$};



%%%%%%%%%%%%%%%%%%%%%%%%%%%%%%%%%%%%%%
%изохора
\begin{scope}[xshift=5cm]
\draw[->] (-.3,0) --++ (0:3cm) node[below,black] {$T$};
\draw[->] (0,-.3) --++ (90:3cm) node[left,black] {$P$};
\node[below left] at (0,0) {$O$};


%red
\draw[dashed] (0,0) --++ (60:.75cm) coordinate (v1);
\draw[thick,red] (v1) --++ (60:2cm) node[black,right] {$V_1$};

\draw[dashed] (0,0) --++ (30:.75cm) coordinate (v1);
\draw[thick,red] (v1) --++ (30:1.5cm) node[black,right] {$V_2>V_1$};

\end{scope}


%%%%%%%%%%%%%%%%%%%%%%%%%%%%%%%%%%%%%%
%изобара
\begin{scope}[xshift=10cm]
\draw[->] (-.3,0) --++ (0:3cm) node[below,black] {$T$};
\draw[->] (0,-.3) --++ (90:3cm) node[left,black] {$V$};
\node[below left] at (0,0) {$O$};


%red
\draw[dashed] (0,0) --++ (60:.75cm) coordinate (v1);
\draw[thick,red] (v1) --++ (60:2cm) node[black,right] {$P_1$};

\draw[dashed] (0,0) --++ (30:.75cm) coordinate (v1);
\draw[thick,red] (v1) --++ (30:1.5cm) node[black,right] {$P_2>P_1$};

\end{scope}


    \end{tikzpicture}
\caption{Графики изопроцессов}
\label{fig:p72}
\end{figure}


% График процесса с $T=\mathrm{const}$ называется изотермой. На рис.\,\ref{fig:p72} (слева) изображены две изотермы газа с разными температурами на $PV$"=диаграмме: \emph{чем выше температура, тем выше лежит соответствующая изотерма на $PV$"=ди\-аграмме}.




 
\end{document}